\clearpage
\par\bigskip\Needspace{9\baselineskip}
\SourceLeaf{222}
\renewcommand{\BookRunningTitle}{Nachlass. Allgemeine Untersuchungen über Kongruenzen}
\pdfbookmark[0]{Allgemeine Untersuchungen über Kongruenzen}{disquisitiones-generales-de-congruentiis}
\begin{center}
{\large ALLGEMEINE UNTERSUCHUNGEN ÜBER KONGRUENZEN.}\par\medskip
LEHRE VON DEN RESTEN, ACHTES KAPITEL.
\end{center}
\BookRule{.18\linewidth}
\renewcommand{\Article}[1]{\par\medskip\Needspace{4\baselineskip}\pdfbookmark[1]{Art. #1}{disquisitiones-generales-art-#1}\centerline{\textbf{#1.}}\nobreak\smallskip}
\Article{330}
Die in den vorigen Abschnitten dargestellten Ergebnisse über Kongruenzen betreffen nur die einfachsten Fälle und wurden meist mit besonderen Methoden gewonnen. In diesem Abschnitt wollen wir versuchen, die Theorie der Kongruenzen, soweit es gegenwärtig möglich ist, auf höhere Grundsätze zurückzuführen, etwa so, wie man die Theorie der \emph{Gleichungen} zu betrachten pflegt. Mit dieser besteht, wie wir schon mehrfach bemerkt haben, eine bemerkenswerte Analogie. Da sich alle algebraischen Kongruenzen mit einer einzigen Unbekannten auf die Form
\[
X\equiv0
\]
bringen lassen, wobei $X$ eine algebraische Funktion der Unbekannten $x$ ohne Brüche ist, müssen zunächst solche Funktionen betrachtet werden.

\Article{331}
Sind $P,Q$ Funktionen der Unbestimmten $x$ von der Form
\[
\begin{gathered}
A+Bx+Cxx+Dx^3+\ldots\\
H+Ix+Kxx+Lx^3+\ldots
\end{gathered}
\]
(solche wollen wir künftig stets einfach als \emph{Funktionen} bezeichnen), und sind in beiden die Koeffizienten gleicher Potenzen von $x$ nach irgendeinem Modul
\SourceLeaf{223}
kongruent, so heißen die \emph{Funktionen nach diesem Modul kongruent}. Offenbar nehmen kongruente Funktionen kongruente Werte an, wenn für die Unbestimmte gleiche oder kongruente Werte eingesetzt werden. Was wir in den Kapiteln \textsc{i} und \textsc{ii} über \emph{Zahlen} bewiesen haben, gilt meist auch für Funktionen. Sind also $P\equiv P'$, $Q\equiv Q'$, $R\equiv R'$ usw., so gelten offenbar $P+Q+R\text{ usw.}\equiv P'+Q'+R'+\text{usw.}$; $P-Q\equiv P'-Q'$; $PQ\equiv P'Q'$; $PQR\text{ usw.}\equiv P'Q'R'\text{ usw.}$ Die Beweise sind sehr leicht und lassen sich ähnlich wie in Kapitel $\mathrm{I}$ führen.

Ist $PQ\equiv R$, so bezeichnen wir die Funktion $Q$ durch $\frac RP$ unter Angabe des Moduls und sagen, $Q$ sei der Quotient, wenn $R$ durch $P$ nach diesem Modul geteilt wird. Offenbar können an die Stelle von $Q$ alle ihr kongruenten Funktionen treten, die wir sämtlich als \emph{einen einzigen} Wert betrachten. Weiter unten werden wir jedoch zeigen, in welchen Fällen ein solcher Quotient mehrere, also inkongruente Werte haben kann.

\Article{332}
Ist der Modul eine Primzahl und enthält der Divisor $Q$ nur einen einzigen Term $Hx^h$, dessen Koeffizient $H$ nicht durch den Modul teilbar ist, ist also nur $H$ nicht $\equiv0$, so kann der Quotient nicht mehrere Werte haben. Denn wären $QA\equiv P$ und $QB\equiv P$, so wäre $Q(A-B)\equiv0$. Es sei nun
\[
Q\equiv\dotstwo+Hx^h+Ix^{h+1}+\text{usw.}
\]
wobei $H$ nicht durch $p$ teilbar sei, und
\[
A-B\equiv Lx^l+Mx^{l+1}+\text{usw.}
\]
wobei $L$ nicht durch $p$ teilbar sei (diese Form wird $A-B$ haben, weil wir voraussetzen, dass $A$ nicht $\equiv B$ ist). Dann wäre $Q(A-B)\equiv HLx^{h+l}+\text{usw.}\equiv0$. Das ist absurd, weil $HL$ nicht $\equiv0$ ist.

Es lassen sich nun leicht Regeln angeben, um die Funktion $P$ durch $Q$ zu teilen, sofern dies möglich ist. Es sei
\[
\begin{aligned}
P&\equiv ax^\alpha+bx^{\alpha+1}+cx^{\alpha+2}+\text{usw.}+kx^\kappa\\
Q&\equiv mx^\mu+nx^{\mu+1}+qx^{\mu+2}+\text{usw.}+tx^\tau
\end{aligned}
\]
wobei $a,k,m,t$ nicht durch den Modul teilbar seien. Dann muss $\alpha$ nicht $<\mu$ und $\kappa$ nicht $<\tau$ sein. Die Division lässt sich ähnlich wie beim gewöhnlichen algebraischen Rechnen ausführen, wenn man nur stets für den Quotienten eine ganze Zahl nimmt. Der Quotient wird nämlich immer
\SourceLeaf{224}
die Form $\frac rm$ haben, was nach dem Modul bestimmt werden muss. Bleibt nun, nachdem $\kappa+\mu-\alpha-\tau+1$ Terme gefunden sind, ein Rest, so hat dieser die Form
\[
Ax^{\kappa+\mu-\tau+1}+Bx^{\kappa+\mu-\tau+2}+\dotstwo+Cx^\kappa
\]
Sind dabei nicht alle Koeffizienten $A,B,C\dotstwo$ $\equiv0$, so lässt sich $P$ nicht durch $Q$ teilen.

Offenbar hätte man die Division auch mit den Termen höchsten Grades, $kx^\kappa$, $tx^\tau$, beginnen können. Die Rechnung wird leichter, wenn man $Q$ auf die Form
\[
mx^\mu(1+qx+rxx+\text{usw.})
\]
bringt. Setzt man $mv\equiv1$, so ergibt sich daraus
\[
\frac PQ\equiv\frac{vP:x^\mu}{1+qx+\text{usw.}}
\]
Danach lässt sich die Division mit den gewöhnlichen Methoden zu Ende führen.

\Article{333}
\textsc{Satz.} \emph{Ist $x\equiv a$ eine Wurzel der Kongruenz $\xi\equiv0$, so lässt sich $\xi$ durch $x-a$ nach dem Modul der Kongruenz teilen.}

\emph{Beweis.} Ließe sie sich nämlich nicht teilen, so wäre $\xi\equiv(x-a)\xi'+b$, wobei $b$ nicht durch den Modul teilbar wäre. Setzt man nun $x$ $\equiv a$, so wird $\xi$ nach Voraussetzung $\equiv0$, also $(x-a)\xi'+b\equiv0$. Dann gilt aber auch $(x-a)\xi'\equiv0$, und daher muss $b$ notwendig $\equiv0$ sein.

\Article{334}
\textsc{Aufgabe.} \emph{Zu zwei gegebenen Funktionen ihren gemeinsamen Teiler (höchsten Grades) nach einem gegebenen Modul finden.}

\emph{Lösung.} Die Funktionen seien $A,B$, und der Grad von $A$ sei mindestens so groß wie der von $B$. Man teile $A$ durch $B$. Geht dies ohne Rest, so ist $B$ der gesuchte gemeinsame Teiler. Bleibt ein Rest $C$, so hat dieser einen niedrigeren Grad als $B$. Es sei also
\[
A\equiv aB+C,\quad B\equiv bC+D,\quad C\equiv cD+E,\quad\text{usw.}
\]
wobei $A,B,C,D,a,b,c$ usw. Funktionen seien und die Grade der Funktionen $A,B,C,D$ usw. eine abnehmende Folge bilden. Geht schließlich eine Division auf, etwa $D\equiv dE$, so ist der letzte Divisor der gesuchte gemeinsame Teiler. Geht keine auf, so gelangt man schließlich zu einem Rest, der keinen Grad hat,
\SourceLeaf{225}
also zu einer \emph{Zahl}. In diesem Fall haben die Funktionen $A,B$ aber keinen gemeinsamen Teiler.

\emph{Beweis.} Teilt der Divisor $E$ die vorhergehende Funktion ohne Rest, so teilt er, wie leicht ersichtlich, alle vorhergehenden. Daher ist $E$ ein gemeinsamer Teiler der Funktionen $A,B$. Damit ist die erste Behauptung bewiesen. Gäbe es aber einen Teiler höheren Grades, etwa $E'$, so würde dieser wegen $C\equiv A-aB$ auch $C$ und ebenso $D$ usw., also auch $E$ teilen: Eine Funktion höheren Grades würde eine Funktion niedrigeren Grades teilen. Das ist absurd. Damit ist die zweite Behauptung bewiesen. Hieraus erkennt man auch, dass man, wenn ein gemeinsamer Teiler irgendeines Grades existiert, nicht zu einem Rest ohne Grad gelangen kann. Sonst würde nämlich eine Funktion ohne Grad durch eine Funktion irgendeines Grades geteilt, was absurd ist.

\Article{335}
\textsc{Satz.} \emph{Sind $A,B$ nach dem Modul $p$ teilerfremde Funktionen und hat $A$ den Grad $\alpha$, $B$ den Grad $\beta$, so lassen sich Funktionen $P,Q$ mit Graden finden, die jeweils $<\beta$, $<\alpha$ sind, sodass}
\[
PA+QB\equiv1\;(\mathit{mod.}\,p)
\]

\Needspace{12\baselineskip}
\emph{Beweis.} In diesem Fall gilt nämlich
\[
A\equiv aB+C,\quad B\equiv bC+D,\quad\text{usw.}\quad K\equiv kL+M
\]
wobei die Grade der Funktionen $A,B,C,D,\dotstwo K,L,M$ fortwährend abnehmen und $M$ keinen Grad hat. Man bilde nun die Folgen
\[
\begin{gathered}
a,\quad a',\quad a'',\quad a''',\ldots a^{(\kappa)}\\
1,\quad b,\quad b',\quad b'',\ldots b^{(\kappa-1)}
\end{gathered}
\]
so, dass
\[
\begin{array}{llll}
a'\equiv ba+1&a''\equiv ca'+a&a'''\equiv da''+a'&\text{usw.}\\
b'\equiv cb+1&b''\equiv db'+b&b'''\equiv eb''+b'&\text{usw.}
\end{array}
\]
Dann gilt
\[
A-aB\equiv+C,\quad bA-a'B\equiv-D,\quad b'A-a''B\equiv+E,\quad\text{usw.}
\]
wie leicht ersichtlich, und daher schließlich
\[
b^{(\kappa-1)}A-a^{(\kappa)}B\equiv\pm M
\]
Es sei nun $\frac1{\pm M}\equiv\mu$. Setzt man $P\equiv\mu b^{(\kappa-1)}$, $Q\equiv-\mu a^{(\kappa)}$, so erhält man
\SourceLeaf{226}
\[
PA+QB\equiv1
\]
Weiter ist offenbar
\[
\begin{aligned}
\text{Grad von }B+\text{Grad von }a&=\text{Grad von }A.\\
\text{Grad von }C+\text{Grad von }b&=\pdot\text{Grad von }B\\
&\text{usw.}\\
\text{Grad von }L+\text{Grad von }k&=\text{Grad von }K
\end{aligned}
\]
Also gilt
\[
\text{Grad von }L+\text{Summe der Grade von }a,b,\dotstwo k=\text{Grad von }A
\]
Der Grad von $a^{(\kappa)}$ ist offenbar und damit auch
\[
\text{Grad von }Q=\text{Summe der Grade von }a,b,c,\dotstwo\text{ d. h. }=\alpha-\text{Grad von }L
\]
Somit ist
\[
\text{Grad von }P=\beta-\text{Grad von }L\qquad\text{Q. E. D.}
\]

\Article{336}
Hieraus folgt, dass man, wenn $M$ der gemeinsame Teiler höchsten Grades der Funktionen $A,B$ ist, stets setzen kann
\[
AP+BQ\equiv M
\]

Beispiele zum vorigen Satz lasse ich der Kürze halber weg. Die Leser werden es aber nicht versäumen, sich durch Beispiele im Umgang mit solchen Aufgaben zu üben. Es lohnt sich auch, darauf hinzuweisen, dass der vorige Satz für Funktionen gilt, die ohne einen Modul betrachtet werden, sofern ihre Koeffizienten rationale Zahlen sind. Das ergibt sich von selbst aus der Art des Beweises. Wir können uns jedoch nicht damit aufhalten. Ähnliches wird der Leser im Folgenden auch ohne Hinweis bemerken.

Hat $A$ weder mit $B$ noch mit $C$ einen gemeinsamen Teiler irgendeines Grades, so hat sie auch mit dem Produkt $BC$ keinen gemeinsamen Teiler. Es sei nämlich
\[
PA+QB\equiv1,\quad\text{dann gilt}\quad PAC+QBC\equiv C
\]
Hätte nun $A$ mit $BC$ einen gemeinsamen Teiler $M$, so würde dieser auch $C$ teilen, entgegen der Voraussetzung. Allgemein ist also eine Funktion $A$, die zu $B,C,D$ usw. teilerfremd ist, auch zum Produkt aller teilerfremd.

Haben $A,B,C,D$ usw. keinen allen gemeinsamen Teiler, so lässt sich erreichen, dass
\SourceLeaf{227}
\[
PA+QB+RC+SD+\text{usw.}\equiv1
\]
Der gemeinsame Teiler höchsten Grades von $A$ und $B$ sei $M$, der von $M$ und $C$ sei $M'$, der von $M'$ und $D$ sei $M''$ usw. Offenbar hat das letzte Glied dieser Folge nach Voraussetzung keinen Grad. Daher lässt sich setzen
\[
aA+bB\equiv M,\quad mM+cC\equiv M',\quad m'M'+dD\equiv M'',\quad\text{usw.}
\]
Nach den entsprechenden Einsetzungen ergibt sich die Richtigkeit des Satzes.

\Article{337}
\textsc{Satz.} \emph{Sind $A,B,C$ usw. nach dem Modul $p$ zueinander teilerfremde Funktionen (haben also je zwei keinen gemeinsamen Teiler), und ist die Funktion $M$ nach demselben Modul durch jede einzelne teilbar, so ist sie auch durch das Produkt aller teilbar.}

\emph{Beweis.} Man kann nämlich $PA+QB\equiv1$ setzen, und daher gilt
\[
\frac MA Q+\frac MB P\equiv\frac M{AB}
\]
Da nun $C$ zu $AB$ teilerfremd ist, ist $M$ auch durch $ABC$ und ebenso durch $ABCD$ usw. teilbar.

\Needspace{14\baselineskip}
\Article{338}
Hat die Kongruenz $\xi\equiv0$ die Wurzeln $x\equiv a$, $x\equiv b$, $x\equiv c$ usw., so lässt sich $\xi$ durch das Produkt aus $(x-a)$, $(x-b)$, $(x-c)$ usw. teilen. Da nämlich $a,b,c\pdot$ usw. als inkongruent vorausgesetzt werden, sind die Funktionen $x-a,x-b,x-c$ usw. zueinander teilerfremd. Da $\xi$ durch jede einzelne teilbar ist, ist sie auch durch das Produkt aller teilbar. Die Anzahl der Wurzeln kann daher den Grad der Kongruenz nicht übersteigen. Das ist der versprochene Beweis dieses Satzes.

Zugleich zeigt sich aber, dass die Lösung der Kongruenzen nur einen Teil einer viel tiefer gehenden Untersuchung bildet, nämlich der Zerlegung von Funktionen in Faktoren. Offenbar hat die Kongruenz $\xi\equiv0$ keine reellen Wurzeln, wenn $\xi$ keine Faktoren ersten Grades besitzt. Das hindert aber nicht daran, dass sich $\xi$ in Faktoren zweiten, dritten oder höheren Grades zerlegen lässt, aus denen man ihr gewissermaßen \emph{imaginäre} Wurzeln zuschreiben könnte. Hätten wir uns tatsächlich eine ähnliche Freiheit nehmen und solche imaginären Größen einführen wollen, wie es die neueren Mathematiker getan haben,
\SourceLeaf{228}
so hätten wir alle folgenden Untersuchungen unvergleichlich verkürzen können. Dennoch haben wir es vorgezogen, alles aus Grundsätzen abzuleiten\begingroup\renewcommand{\thefootnote}{*)}\footnote{Vielleicht werden wir bei einer anderen Gelegenheit unsere Ansicht hierzu ausführlicher darlegen.}\endgroup.

\Article{339}
\emph{Funktionen} heißen nach einem bestimmten Modul \emph{prim}, wenn sie sich nach diesem Modul durch keine Funktionen niedrigeren Grades teilen lassen.

So sind alle Funktionen ersten Grades prim; Funktionen zweiten Grades sind entweder prim oder aus zwei Funktionen ersten Grades zusammengesetzt. Daher ist $\xi$ eine Primfunktion zweiten Grades, wenn die Kongruenz $\xi\equiv0$ keine reellen Wurzeln zulässt. Zum Beispiel ist $xx+x+1$ für den Modul $5$ prim, weil
\[
xx+x+1\equiv(x-2)^2-3\;(\mathrm{mod.}\,5)
\]
und $3$ ein quadratischer Nichtrest von $5$ ist.

Diese Primfunktionen verdienen vor allem unsere Aufmerksamkeit. Zwar können andere als solche ersten Grades nicht zum Auffinden reeller Wurzeln dienen; eine eingehendere Betrachtung empfiehlt sich jedoch sowohl wegen ihrer bemerkenswerten Eigenschaften als auch wegen anderer vorzüglicher Wahrheiten, die sich aus ihnen ableiten lassen.

\Needspace{20\baselineskip}
\Article{340}
\textsc{Satz.} \emph{Jede Funktion ist entweder prim oder aus Primfunktionen zusammengesetzt. Im zweiten Fall lässt sie sich nur auf eine einzige Weise aus Primfunktionen zusammensetzen.}

\emph{Beweis.} Ist nämlich die vorgelegte Funktion $A$ nicht prim, so ist sie durch eine andere Funktion $B$ niedrigeren Grades teilbar. Ist $B$ keine Primfunktion, so ist sie durch eine andere Funktion $C$ niedrigeren Grades teilbar. Fährt man so fort, so muss man schließlich zu einer Primfunktion gelangen. Sonst wäre diese Folge unendlich, was absurd ist, da die Grade fortwährend abnehmen. Ist nun die letzte Primfunktion $L$, so teilt sie alle vorhergehenden. Daher ist $A\equiv LA'$, und $A'$ hat einen niedrigeren Grad als $A$. Wiederholt man dies, so erhält man $A'\equiv L'A''$ usw. Offenbar muss man schließlich zu einer Primfunktion gelangen, und daher ist $A$ $\equiv$ einem Produkt aus Primfunktionen $L,L',L''$ usw. $\pdot$ Damit ist die erste Behauptung bewiesen.

Wäre nun auch $A\equiv MM'M''$ usw., und wären nicht alle $L,L',L''$ usw. dieselben wie alle $M,M',M''$ usw., so lasse man die Funktionen weg, die beiden Folgen gemeinsam sind.
\SourceLeaf{229}
Es mögen $\lambda,\lambda',\lambda''\dotstwo$; $\mu,\mu',\mu'',\ldots$ übrig bleiben. Dann ist $\mu$ zu $\lambda,\lambda',\lambda''$ usw., also auch zum Produkt $\lambda\lambda'\lambda''$ usw. teilerfremd. Dennoch müsste gelten
\[
\lambda\lambda'\lambda''\dotstwo\equiv\mu\mu'\mu''\dotstwo\text{ d. h. }\frac{\lambda\lambda'\lambda''\dotstwo}{\mu}\equiv\mu'\mu''\dotstwo\quad\text{ein Widerspruch}
\]
